\section*{Chapter \ref{ch:g9:earth-resources} --- Resources from the Earth}

\begin{solution}{exo:g9:earth-resources:1}
(a) \omterm{def:g9:earth-resources:aggregate}{Aggregate}. (b) \omterm{def:g9:earth-resources:reserve}{Reserve}. (c) \omterm{def:g9:earth-resources:ore-deposit}{Ore deposit}. (d) \omterm{def:g9:earth-resources:ore-deposit}{Concentration factor}.
\end{solution}

\begin{solution}{exo:g9:earth-resources:2}
(a) In \omterm{def:g1:landscapes-around-us:water}{river} \omterm{def:g3:stones-pebbles-sand:pebble}{gravels}. (b) In cracks, from hot water. (c) In the \omterm{def:g7:volcanoes:magma}{magma}.
(d) This is the first step of a \omterm{def:g3:stones-pebbles-sand:pebble}{river-gravel} deposit: the freed gold is
then carried by the \omterm{def:g1:landscapes-around-us:water}{river}.
\end{solution}

\begin{solution}{exo:g9:earth-resources:3}
(b), (d), (a), (e), (c).
\end{solution}

\begin{solution}{exo:g9:earth-resources:4}
(a) Not renewable: it takes millions of years to form. (b) Renewable: the
\omterm{def:g4:clouds-rain-snow:precipitation}{rain} refills what is taken. (c) Not renewable: an \omterm{def:g9:earth-resources:ore-deposit}{ore deposit} forms over
very long times. (d) Not renewable at that rate: the \omterm{def:g6:water-underground:water-table}{water table} falls and
the \omterm{def:g6:water-underground:spring}{wells} end up dry.
\end{solution}

\begin{solution}{exo:g9:earth-resources:5}
(a) $4 \div 0.007 \approx 571$, about $10^3$ (between $10^2$ and $10^3$,
closer to $10^3$). (b) $0.5 \div 0.0002 = 2500$, about $10^3$. (c) $20
\div 5.6 \approx 3.6$, about 4: between $10^0$ and $10^1$.
(d) Iron is already very abundant in the \omterm{def:g6:inside-the-earth:layers}{crust}: an ore only a few times
richer than ordinary \omterm{def:g3:stones-pebbles-sand:rock}{rock} is enough.
\end{solution}

\begin{solution}{exo:g9:earth-resources:6}
(a) $19.3 \div 2.5 \approx 7.7$ times. (b) Gold: $19.3 \times 0.01 =
\qty{0.193}{g}$; \omterm{def:g3:stones-pebbles-sand:rock}{rock}: $2.5 \times 0.01 = \qty{0.025}{g}$. (c) When the
whole bed moves, the heavy gold grains sink between the lighter \omterm{def:g3:stones-pebbles-sand:sand}{sand} and
\omterm{def:g3:stones-pebbles-sand:pebble}{pebbles}, down to the solid \omterm{def:g3:stones-pebbles-sand:rock}{bedrock}, where they stop.
\end{solution}

\begin{solution}{exo:g9:earth-resources:7}
(a) Gold, tin and chromium. (b) Between $10^1$ and $10^2$, close to
$10^2$. (c) $3000 \div 3.7 \approx 800$ times. (d) Chromium (or tin, or
gold): the higher the \omterm{def:g9:earth-resources:ore-deposit}{concentration factor}, the rarer the places where
nature has concentrated the metal enough.
\end{solution}

\begin{solution}{exo:g9:earth-resources:8}
(a) $18 \times 365 = \qty{6570}{kg}$, about \qty{6.6}{t}. (b) $5 \times
10^{10}$ tonnes. (c) $25000 \div 18 \approx 1389$ days, almost four years.
(d) The \omterm{def:g3:stones-pebbles-sand:sand}{sand} dug from the \omterm{def:g1:landscapes-around-us:water}{river} no longer travels downstream to the \omterm{def:g1:landscapes-around-us:water}{coast};
the waves keep carrying \omterm{def:g3:stones-pebbles-sand:sand}{sand} away from the beach, but less arrives to
replace it.
\end{solution}

\begin{solution}{exo:g9:earth-resources:9}
(a) From tiny plants and animals of the \omterm{def:g2:water-on-earth:ocean}{sea}, such as diatoms, buried in the
\omterm{def:g3:stones-pebbles-sand:sand}{mud} of the \omterm{def:g2:water-on-earth:ocean}{sea} floor. (b) Between about 60 and \qty{120}{\degreeCelsius}.
(c) They are lighter than water, which fills the \omterm{def:g6:water-underground:porosity}{pores} of the \omterm{def:g3:stones-pebbles-sand:rock}{rocks}.
(d) An \omterm{def:g6:water-underground:permeability}{impermeable} layer, such as shale; they gather under it, often at
the top of an upfold, gas above oil.
\end{solution}

\begin{solution}{exo:g9:earth-resources:10}
(a) $2.9 \times 10^7$ tonnes. (b) $2.9 \times 10^7 \div 100 = 2.9 \times
10^5$ tonnes a year, roughly. (c) The ore is 100 times heavier than the
copper it holds: $2.9 \times 10^7 \times 100 = 2.9 \times 10^9$ tonnes.
(d) Every tonne of ore gives only \qty{10}{kg} of copper, so billions of
tonnes of \omterm{def:g3:stones-pebbles-sand:rock}{rock} have to be dug.
\end{solution}

\begin{solution}{exo:g9:earth-resources:11}
(a) At 23~million tonnes a year it would have lasted about $280 \div 23
\approx 12$ years, until the early 1980s; yet 697~million tonnes were mined,
much more than that \omterm{def:g9:earth-resources:reserve}{reserve}. (b) $980 \div 280 = 3.5$ times. (c) $980 \div
23 \approx 43$ years. (d) For example: new deposits are discovered; better
techniques or higher prices make poorer deposits worth mining.
\end{solution}

\begin{solution}{exo:g9:earth-resources:12}
(a) $5 \times 10^1$ years. (b) $2 \times 10^3 \div 50 = 40$ times. (c) $3
\times 10^8 \div (5 \times 10^1) = 6 \times 10^6$ times. (d) The coal took
millions of times longer to form than we take to use it: nature does not
renew it on our time scale.
\end{solution}

\begin{solution}{pb:g9:earth-resources:1}
\textbf{1.} $100 \times 1000 = \num{100000}$~kg $= 1 \times 10^5$~kg.

\textbf{2.} $2000 \times 300 = \num{600000}$~t $= 6 \times 10^5$~t.

\textbf{3.} $6 \times 10^5 \div 10 = 6 \times 10^4$~t per year.

\textbf{4.} $0.40 \times 6 \times 10^4 = 2.4 \times 10^4$~t of \omterm{def:g3:stones-pebbles-sand:sand}{sand} per
year.

\textbf{5.} $1200 \times 500 \times 5 = 3 \times 10^6$~m$^3$.

\textbf{6.} Two thirds of it: $2 \times 10^6$~m$^3$.

\textbf{7.} $2 \times 10^6 \times 1.6 = 3.2 \times 10^6$~t.

\textbf{8.} $\num{24000} + \num{136000} = \num{160000}$~t $= 1.6 \times
10^5$~t per year.

\textbf{9.} $\num{160000} \div 2000 = 80$ times more.

\textbf{10.} No: the pit is emptied 80 times faster than the \omterm{def:g1:landscapes-around-us:water}{river}
renews the \omterm{def:g3:stones-pebbles-sand:sand}{sand}. On the scale of the town, this \omterm{def:g3:stones-pebbles-sand:sand}{sand} is a finite resource.

\textbf{11.} $\num{24000} \div 4 = \num{6000}$~t of \omterm{def:g3:stones-pebbles-sand:sand}{sand} a year.

\textbf{12.} $3.2 \times 10^6 \div (1.6 \times 10^5) = 20$: \textbf{20 years}.
\end{solution}
